%! Author = oliver
%! Date = 12/02/2026

\title{Valg af DB}
% Document
\section{Valg af DB}

    Til valg af vores database er der to valg. SQL eller NoSQL. Her er der mange forskellige implementationer af disse,
    hvor de mest kendte er MongoDB til NoSQL og PostgreSQL til SQL, og det er dem vi vil vælge imellem. Der er fordele og
    ulemper med begge.

    \subsection{MongoDB}

    MongoDB er en NoSQL dokumentdatabase, som lagrer data i JSON-lignende dokumenter (BSON).
    I stedet for tabeller og relationer arbejder man med collections og indlejrede dokumenter.
    Det betyder, at datastrukturen er fleksibel og ikke kræver et fast defineret skema.

    En fordel ved MongoDB er, at den er velegnet til applikationer med hurtigt skiftende
    datastrukturer eller store mængder ustruktureret data. Den kan være nem at skalere
    horisontalt og giver høj performance ved simple læse- og skriveoperationer.

    Ulempen ved MongoDB i denne kontekst er manglen på stærke relationer mellem data.
    I en applikation med mange relationer – såsom spillere, hold, turneringer, kampe og resultater –
    kan det blive komplekst at håndtere relationer og sikre dataintegritet.
    Der findes ikke samme niveau af indbygget referentiel integritet som i relationelle databaser,
    hvilket kan føre til inkonsistent data, hvis det ikke håndteres korrekt i applikationslaget.

    \subsection{PostgreSQL}

    PostgreSQL er en relationel SQL-database, som organiserer data i tabeller med
    definerede relationer mellem dem \cite{postgres_docs2026}. Den understøtter ACID-egenskaber, fremmednøgler,
    transaktioner og avancerede forespørgsler via SQL.

    Fordelen ved PostgreSQL i denne webapplikation er den stærke håndtering af relationer.
    En kajakpolo-applikation indeholder mange naturlige relationer, eksempelvis:

    \begin{itemize}
        \item En spiller tilhører et hold
        \item Et hold deltager i en turnering
        \item En turnering består af flere kampe
        \item En kamp har to hold og et resultat
        \item Brugere kan have forskellige roller (spiller, arrangør, administrator)
    \end{itemize}

    Disse relationer kan modelleres direkte med primær- og fremmednøgler,
    hvilket sikrer dataintegritet og konsistens. PostgreSQL understøtter desuden
    komplekse joins, hvilket gør det effektivt at hente relateret data, f.eks.
    alle kampe i en turnering eller alle spillere på et bestemt hold.

    Da systemet kræver struktureret data med klare relationer og høj konsistens,
    er en relationel database det mest hensigtsmæssige valg.

    \subsection{Konklusion}

    På baggrund af systemets krav om stærke relationer, dataintegritet og
    struktureret data vælges PostgreSQL som database.
    En relationel model passer naturligt til domænet med hold, spillere,
    turneringer og kampe, og PostgreSQL giver robuste værktøjer til at
    håndtere disse relationer effektivt.
